
\documentclass[Main.tex]{subfiles}
\begin{document}
%%%Phần mở đầu
\chapter{Tốc Độ Phản Ứng Hóa Học}
\section{Tốc độ phản ứng hóa học}
\begin{Muctieu}
    \begin{itemize}
        \item Nêu được khái niệm tốc độ phản ứng và tốc độ trung bình của phản ứng.
        \item Trình bày được các yếu tố ảnh hưởng đến tốc độ phản ứng (nồng độ, nhiệt độ, áp suất, diện tích bề mặt, chất xúc tác).
        \item Giải thích được ảnh hưởng của các yếu tố đến tốc độ phản ứng dựa trên thuyết va chạm.
        \item Vận dụng được các yếu tố ảnh hưởng để điều khiển tốc độ phản ứng trong thực tiễn.
    \end{itemize}
\end{Muctieu}
\begin{kd}
	\immini{Trong cuộc sống hàng ngày, chúng ta thường thấy có những phản ứng hóa học diễn ra rất nhanh, ví dụ như sự cháy của xăng, gas, hay sự nổ của bom. Ngược lại, có những phản ứng diễn ra rất chậm, ví dụ như sự gỉ sét của kim loại, hay sự lên men của rượu. Vậy, điều gì quyết định tốc độ của một phản ứng hóa học? Và làm thế nào để chúng ta có thể làm cho một phản ứng diễn ra nhanh hơn hoặc chậm hơn?}{}%\includegraphics[height=4.5cm]{Images/anhhoahoc10/TOC_DO_PU/chay.png}\hfill\includegraphics[height=4.5cm]{Images/anhhoahoc10/TOC_DO_PU/gi_sat.png}
	Câu trả lời nằm ở khái niệm về tốc độ phản ứng hóa học và các yếu tố ảnh hưởng đến nó. Trong bài học này, chúng ta sẽ cùng nhau tìm hiểu về tốc độ phản ứng, cách đo lường nó, cũng như các yếu tố có thể làm thay đổi tốc độ của một phản ứng hóa học.
\end{kd}
%%%Phần lý thuyết
\subsection{Nội dung bài học}
\subsubsection{Khái niệm tốc độ phản ứng}
    \Noibat[\maunhan][][][]{Tốc độ phản ứng}
    \begin{ghinho}
        \indam[\maunhan]{Tốc độ phản ứng} là độ biến thiên nồng độ của một chất phản ứng hoặc sản phẩm trong một đơn vị thời gian.
    \end{ghinho}
    \Noibat[\maunhan][][][]{Tốc độ trung bình của phản ứng}
    \begin{ghinho}
        \indam[\maunhan]{Tốc độ trung bình của phản ứng} được tính bằng độ biến thiên nồng độ của một chất phản ứng hoặc sản phẩm chia cho khoảng thời gian xảy ra sự biến thiên đó.
        \[
            v_{tb} = \frac{\Delta C}{\Delta t}
        \]
        Trong đó:
        \begin{itemize}
            \item $\Delta C$ là độ biến thiên nồng độ.
            \item $\Delta t$ là khoảng thời gian xảy ra sự biến thiên nồng độ.
        \end{itemize}
    \end{ghinho}
\subsubsection{Các yếu tố ảnh hưởng đến tốc độ phản ứng}
    \Noibat[\maunhan][][][]{Nồng độ}
    \begin{ghinho}
        Khi tăng nồng độ của các chất phản ứng, tốc độ phản ứng thường tăng.
    \end{ghinho}
    \Noibat[\maunhan][][][]{Nhiệt độ}
    \begin{ghinho}
        Khi tăng nhiệt độ, tốc độ phản ứng thường tăng.
    \end{ghinho}
    \Noibat[\maunhan][][][]{Áp suất}
    \begin{ghinho}
        Đối với các phản ứng có chất khí tham gia, khi tăng áp suất, tốc độ phản ứng thường tăng.
    \end{ghinho}
    \Noibat[\maunhan][][][]{Diện tích bề mặt}
    \begin{ghinho}
        Đối với các phản ứng có chất rắn tham gia, khi tăng diện tích bề mặt tiếp xúc, tốc độ phản ứng thường tăng.
    \end{ghinho}
    \Noibat[\maunhan][][][]{Chất xúc tác}
    \begin{ghinho}
        \indam{Chất xúc tác} là chất làm tăng tốc độ phản ứng nhưng không bị tiêu thụ trong quá trình phản ứng.
    \end{ghinho}
\subsubsection{Thuyết va chạm}
    \Noibat[\maunhan][][][]{Thuyết va chạm}
    \begin{tomtat}
        \begin{itemize}
            \item Các phân tử chất phản ứng phải va chạm với nhau thì phản ứng mới xảy ra.
            \item Không phải tất cả các va chạm đều dẫn đến phản ứng, mà chỉ có các va chạm có đủ năng lượng (năng lượng hoạt hóa) và đúng hướng mới tạo thành sản phẩm.
        \end{itemize}
    \end{tomtat}
    \Noibat[\maunhan][][][]{Giải thích ảnh hưởng của các yếu tố}
    \begin{hopdongian}
        Các yếu tố như nồng độ, nhiệt độ, áp suất, diện tích bề mặt, chất xúc tác đều làm tăng số va chạm hiệu quả giữa các phân tử chất phản ứng, từ đó làm tăng tốc độ phản ứng.
    \end{hopdongian}
%%%Phần bài tập
\subsection{Các dạng bài tập}
\begin{dang}{Dạng 1: Tính tốc độ trung bình của phản ứng}
\end{dang}
\begin{phuongphap}
    Sử dụng công thức: $v_{tb} = \frac{\Delta C}{\Delta t}$
\end{phuongphap}
\Noibat[\maunhan][][\faBookmark][]{Ví dụ mẫu}
\begin{vd}
    Cho phản ứng: $2SO_2(g) + O_2(g) \rightarrow 2SO_3(g)$. Nồng độ $SO_2$ giảm từ 0.5M xuống 0.2M trong 5 phút. Tính tốc độ trung bình của phản ứng trong khoảng thời gian này.
	\loigiai{
        \[
            v_{tb} = \frac{0.5 - 0.2}{5} = 0.06 \text{ M/phút}
        \]
    }
\end{vd}
%%%%%=====================Bài tập tự luyện Dạng 1==========================%%%
\Noibat[\maunhan][][\faBook][]{Bài tập tự luyện}

\phan{Bài tập tự luận và trả lời ngắn}
%%%=============SOẠN BT===============%%%
\Opensolutionfile{ansbth}[Ans/LGBT-C06B01_Dang1]
\Opensolutionfile{ansbt}[Ans/AnsBT-C06B01_Dang1]
    %%%%%============BT_1================%%%%%%
    \begin{bt}
        Cho phản ứng: $H_2(g) + I_2(g) \rightarrow 2HI(g)$. Nồng độ $H_2$ giảm từ 0.6M xuống 0.3M trong 10 phút. Tính tốc độ trung bình của phản ứng trong khoảng thời gian này.
        \loigiai{
            \[
                v_{tb} = \frac{0.6 - 0.3}{10} = 0.03 \text{ M/phút}
            \]
        }
    \end{bt}
    %%%%%============BT_2================%%%%%%
     \begin{bt}
        Cho phản ứng: $2NO(g) + O_2(g) \rightarrow 2NO_2(g)$. Nồng độ $NO$ giảm từ 0.8M xuống 0.4M trong 8 phút. Tính tốc độ trung bình của phản ứng trong khoảng thời gian này.
        \loigiai{
            \[
                v_{tb} = \frac{0.8 - 0.4}{8} = 0.05 \text{ M/phút}
            \]
        }
    \end{bt}
    %%%%%============BT_3================%%%%%%
    \begin{bt}
        Cho phản ứng: $2N_2O_5(g) \rightarrow 4NO_2(g) + O_2(g)$. Nồng độ $N_2O_5$ giảm từ 1.0M xuống 0.2M trong 20 phút. Tính tốc độ trung bình của phản ứng trong khoảng thời gian này.
        \loigiai{
            \[
                v_{tb} = \frac{1.0 - 0.2}{20} = 0.04 \text{ M/phút}
            \]
        }
    \end{bt}
    %%%%%============BT_4================%%%%%%
    \begin{bt}
        Cho phản ứng: $CaCO_3(s) \rightarrow CaO(s) + CO_2(g)$. Lượng $CaCO_3$ giảm từ 20g xuống 10g trong 10 phút. Tính tốc độ trung bình của phản ứng theo khối lượng $CaCO_3$.
        \loigiai{
            \[
                v_{tb} = \frac{20 - 10}{10} = 1 \text{ g/phút}
            \]
        }
    \end{bt}
     %%%%%============BT_5================%%%%%%
    \begin{bt}
        Cho phản ứng: $Zn(s) + 2HCl(aq) \rightarrow ZnCl_2(aq) + H_2(g)$. Lượng $Zn$ giảm từ 13g xuống 6.5g trong 5 phút. Tính tốc độ trung bình của phản ứng theo khối lượng $Zn$.
        \loigiai{
            \[
                v_{tb} = \frac{13 - 6.5}{5} = 1.3 \text{ g/phút}
            \]
        }
    \end{bt}
     %%%%%============BT_6================%%%%%%
    \begin{bt}
        Cho phản ứng: $2KMnO_4(aq) + 10FeSO_4(aq) + 8H_2SO_4(aq) \rightarrow K_2SO_4(aq) + 2MnSO_4(aq) + 5Fe_2(SO_4)_3(aq) + 8H_2O(l)$. Nồng độ $KMnO_4$ giảm từ 0.1M xuống 0.02M trong 16 phút. Tính tốc độ trung bình của phản ứng trong khoảng thời gian này.
        \loigiai{
            \[
                v_{tb} = \frac{0.1 - 0.02}{16} = 0.005 \text{ M/phút}
            \]
        }
    \end{bt}
    %%%%%============BTS_1================%%%%%%
    \begin{bt}
        Nêu công thức tính tốc độ trung bình của phản ứng.
        \shortans{$v_{tb} = \frac{\Delta C}{\Delta t}$}
        \loigiai{
            Công thức tính tốc độ trung bình của phản ứng là:
            \[
                 v_{tb} = \frac{\Delta C}{\Delta t}
            \]
        }
    \end{bt}
     %%%%%============BTS_2================%%%%%%
    \begin{bt}
        Nêu ý nghĩa của tốc độ trung bình của phản ứng.
        \shortans{Tốc độ trung bình của phản ứng cho biết mức độ nhanh hay chậm của phản ứng trong một khoảng thời gian xác định.}
        \loigiai{
            Tốc độ trung bình của phản ứng cho biết mức độ nhanh hay chậm của phản ứng trong một khoảng thời gian xác định.
        }
    \end{bt}
    %%%%%============BTS_3================%%%%%%
     \begin{bt}
        Đơn vị thường dùng của tốc độ phản ứng là gì?
        \shortans{mol/l.s hoặc M/s}
        \loigiai{
            Đơn vị thường dùng của tốc độ phản ứng là mol/l.s hoặc M/s.
        }
    \end{bt}
    %%%%%============BTS_4================%%%%%%
    \begin{bt}
        Tốc độ phản ứng có thể âm không?
        \shortans{Không}
        \loigiai{
            Tốc độ phản ứng là đại lượng luôn dương.
        }
    \end{bt}
    %%%%%============BTS_5================%%%%%%
    \begin{bt}
        Tốc độ phản ứng là gì?
        \shortans{Tốc độ phản ứng là độ biến thiên nồng độ của một chất phản ứng hoặc sản phẩm trong một đơn vị thời gian.}
        \loigiai{
            Tốc độ phản ứng là độ biến thiên nồng độ của một chất phản ứng hoặc sản phẩm trong một đơn vị thời gian.
        }
    \end{bt}
\Closesolutionfile{ansbt}
\Closesolutionfile{ansbth}
%\bangdapanSA{AnsBT-C06B01_Dang1}
%%%%============Phần trắc nghiệm============%%%
\phan{Trắc nghiệm nhiều lựa chọn}
%%%=============SOẠN EX===============%%%
\Opensolutionfile{ansex}[Ans/LGEX-C06B01_Dang1]
\Opensolutionfile{ans}[Ans/Ans-C06B01_Dang1]
    %%%%%============EX_1================%%%%%%
    \begin{ex}
        Cho phản ứng: $2NO(g) + O_2(g) \rightarrow 2NO_2(g)$. Nồng độ $NO$ giảm từ 0.6M xuống 0.2M trong 10 phút. Tính tốc độ trung bình của phản ứng theo $NO$.
        \choice
        {0.02 M/phút}
        {0.03 M/phút}
        {\True 0.04 M/phút}
        {0.06 M/phút}
        \loigiai{
            \[
                v_{tb} = \frac{0.6 - 0.2}{10} = 0.04 \text{ M/phút}
            \]
        }
    \end{ex}
    %%%%%============EX_2================%%%%%%
    \begin{ex}
        Cho phản ứng: $2H_2O_2(aq) \rightarrow 2H_2O(l) + O_2(g)$. Nồng độ $H_2O_2$ giảm từ 1.0M xuống 0.5M trong 20 phút. Tính tốc độ trung bình của phản ứng theo $H_2O_2$.
        \choice
        {0.015 M/phút}
        {\True 0.025 M/phút}
        {0.03 M/phút}
        {0.04 M/phút}
        \loigiai{
            \[
                v_{tb} = \frac{1.0 - 0.5}{20} = 0.025 \text{ M/phút}
            \]
        }
    \end{ex}
     %%%%%============EX_3================%%%%%%
    \begin{ex}
        Cho phản ứng: $2SO_2(g) + O_2(g) \rightarrow 2SO_3(g)$. Nồng độ $SO_2$ giảm từ 0.9M xuống 0.3M trong 15 phút. Tính tốc độ trung bình của phản ứng theo $SO_2$.
        \choice
        {0.02 M/phút}
        {\True 0.04 M/phút}
        {0.06 M/phút}
        {0.08 M/phút}
        \loigiai{
            \[
                v_{tb} = \frac{0.9 - 0.3}{15} = 0.04 \text{ M/phút}
            \]
        }
    \end{ex}
    %%%%%============EX_4================%%%%%%
    \begin{ex}
        Cho phản ứng: $CaCO_3(s) \rightarrow CaO(s) + CO_2(g)$. Lượng $CaCO_3$ giảm từ 15g xuống 5g trong 10 phút. Tính tốc độ trung bình của phản ứng theo khối lượng $CaCO_3$.
        \choice
        {0.5 g/phút}
        {\True 1 g/phút}
        {1.5 g/phút}
        {2 g/phút}
        \loigiai{
            \[
                v_{tb} = \frac{15 - 5}{10} = 1 \text{ g/phút}
            \]
        }
    \end{ex}
     %%%%%============EX_5================%%%%%%
    \begin{ex}
        Cho phản ứng: $Zn(s) + 2HCl(aq) \rightarrow ZnCl_2(aq) + H_2(g)$. Lượng $Zn$ giảm từ 15.6g xuống 7.8g trong 13 phút. Tính tốc độ trung bình của phản ứng theo khối lượng $Zn$.
        \choice
        {0.4 g/phút}
        {0.5 g/phút}
        {\True 0.6 g/phút}
        {0.8 g/phút}
        \loigiai{
            \[
                v_{tb} = \frac{15.6 - 7.8}{13} = 0.6 \text{ g/phút}
            \]
        }
    \end{ex}
     %%%%%============EX_6================%%%%%%
    \begin{ex}
       Cho phản ứng: $2KMnO_4(aq) + 10FeSO_4(aq) + 8H_2SO_4(aq) \rightarrow K_2SO_4(aq) + 2MnSO_4(aq) + 5Fe_2(SO_4)_3(aq) + 8H_2O(l)$. Nồng độ $KMnO_4$ giảm từ 0.2M xuống 0.08M trong 12 phút. Tính tốc độ trung bình của phản ứng trong khoảng thời gian này.
        \choice
        {0.005 M/phút}
        {0.008 M/phút}
        {\True 0.01 M/phút}
        {0.012 M/phút}
        \loigiai{
            \[
                v_{tb} = \frac{0.2 - 0.08}{12} = 0.01 \text{ M/phút}
            \]
        }
    \end{ex}
     %%%%%============EX_7================%%%%%%
    \begin{ex}
         Tốc độ phản ứng được đo bằng sự thay đổi của đại lượng nào sau đây?
        \choice
        {Khối lượng}
        {Thể tích}
        {\True Nồng độ}
        {Áp suất}
        \loigiai{Tốc độ phản ứng được đo bằng sự thay đổi nồng độ của chất phản ứng hoặc sản phẩm.}
    \end{ex}
     %%%%%============EX_8================%%%%%%
    \begin{ex}
        Tốc độ phản ứng trung bình được tính bằng:
        \choice
        {Tổng nồng độ các chất phản ứng}
        {Tổng nồng độ các sản phẩm}
        {\True Độ biến thiên nồng độ chia cho thời gian}
        {Độ biến thiên thời gian chia cho nồng độ}
        \loigiai{Tốc độ trung bình của phản ứng được tính bằng độ biến thiên nồng độ chia cho thời gian.}
    \end{ex}
     %%%%%============EX_9================%%%%%%
     \begin{ex}
        Trong phản ứng $A \rightarrow B$, nồng độ chất A giảm từ 1M xuống 0.5M trong 10 phút. Tốc độ trung bình của phản ứng là:
        \choice
        {0.02 M/phút}
        {0.03 M/phút}
        {\True 0.05 M/phút}
        {0.1 M/phút}
        \loigiai{
            \[
                v_{tb} = \frac{1 - 0.5}{10} = 0.05 \text{ M/phút}
            \]
        }
    \end{ex}
     %%%%%============EX_10================%%%%%%
     \begin{ex}
        Trong phản ứng $2A \rightarrow B$, nồng độ chất A giảm từ 1.2M xuống 0.6M trong 20 phút. Tốc độ trung bình của phản ứng là:
        \choice
        {0.01 M/phút}
        {\True 0.03 M/phút}
        {0.04 M/phút}
        {0.06 M/phút}
        \loigiai{
            \[
                v_{tb} = \frac{1.2 - 0.6}{20} = 0.03 \text{ M/phút}
            \]
        }
    \end{ex}
     %%%%%============EX_11================%%%%%%
    \begin{ex}
        Trong phản ứng $A + B \rightarrow C$, nồng độ chất A giảm từ 0.8M xuống 0.4M trong 10 phút. Tốc độ trung bình của phản ứng là:
        \choice
        {0.02 M/phút}
        {0.03 M/phút}
        {\True 0.04 M/phút}
        {0.06 M/phút}
        \loigiai{
            \[
                v_{tb} = \frac{0.8 - 0.4}{10} = 0.04 \text{ M/phút}
            \]
        }
    \end{ex}
     %%%%%============EX_12================%%%%%%
    \begin{ex}
        Trong phản ứng $A \rightarrow 2B$, nồng độ chất A giảm từ 0.5M xuống 0.1M trong 8 phút. Tốc độ trung bình của phản ứng là:
         \choice
        {0.02 M/phút}
        {0.03 M/phút}
        {\True 0.05 M/phút}
        {0.06 M/phút}
        \loigiai{
            \[
                v_{tb} = \frac{0.5 - 0.1}{8} = 0.05 \text{ M/phút}
            \]
        }
    \end{ex}
     %%%%%============EX_13================%%%%%%
    \begin{ex}
        Trong phản ứng $A + B \rightarrow C + D$, nồng độ chất B giảm từ 1.2M xuống 0.6M trong 12 phút. Tốc độ trung bình của phản ứng là:
        \choice
        {0.02 M/phút}
        {\True 0.05 M/phút}
        {0.06 M/phút}
        {0.1 M/phút}
        \loigiai{
            \[
                v_{tb} = \frac{1.2 - 0.6}{12} = 0.05 \text{ M/phút}
            \]
        }
    \end{ex}
     %%%%%============EX_14================%%%%%%
    \begin{ex}
        Trong phản ứng $2A \rightarrow B + C$, nồng độ chất A giảm từ 1.0M xuống 0.4M trong 15 phút. Tốc độ trung bình của phản ứng là:
        \choice
        {0.02 M/phút}
        {\True 0.04 M/phút}
        {0.06 M/phút}
        {0.08 M/phút}
        \loigiai{
            \[
                v_{tb} = \frac{1.0 - 0.4}{15} = 0.04 \text{ M/phút}
            \]
        }
    \end{ex}
     %%%%%============EX_15================%%%%%%
    \begin{ex}
        Trong phản ứng $A + 2B \rightarrow C$, nồng độ chất B giảm từ 0.8M xuống 0.2M trong 10 phút. Tốc độ trung bình của phản ứng là:
        \choice
        {0.04 M/phút}
        {0.05 M/phút}
        {\True 0.06 M/phút}
        {0.08 M/phút}
        \loigiai{
            \[
                v_{tb} = \frac{0.8 - 0.2}{10} = 0.06 \text{ M/phút}
            \]
        }
    \end{ex}
     %%%%%============EX_16================%%%%%%
    \begin{ex}
        Trong phản ứng $A + B \rightarrow 2C$, nồng độ chất A giảm từ 0.7M xuống 0.1M trong 12 phút. Tốc độ trung bình của phản ứng là:
        \choice
         {0.02 M/phút}
        {0.03 M/phút}
        {\True 0.05 M/phút}
        {0.06 M/phút}
        \loigiai{
            \[
                v_{tb} = \frac{0.7 - 0.1}{12} = 0.05 \text{ M/phút}
            \]
        }
    \end{ex}
     %%%%%============EX_17================%%%%%%
    \begin{ex}
        Trong phản ứng $3A \rightarrow B$, nồng độ chất A giảm từ 1.5M xuống 0.3M trong 20 phút. Tốc độ trung bình của phản ứng là:
        \choice
        {0.02 M/phút}
        {0.03 M/phút}
        {\True 0.06 M/phút}
        {0.08 M/phút}
        \loigiai{
            \[
                v_{tb} = \frac{1.5 - 0.3}{20} = 0.06 \text{ M/phút}
            \]
        }
    \end{ex}
     %%%%%============EX_18================%%%%%%
    \begin{ex}
        Trong phản ứng $A + 2B \rightarrow C + D$, nồng độ chất B giảm từ 0.9M xuống 0.3M trong 15 phút. Tốc độ trung bình của phản ứng là:
        \choice
        {0.02 M/phút}
        {0.03 M/phút}
        {\True 0.04 M/phút}
        {0.06 M/phút}
        \loigiai{
            \[
                v_{tb} = \frac{0.9 - 0.3}{15} = 0.04 \text{ M/phút}
            \]
        }
    \end{ex}
     %%%%%============EX_19================%%%%%%
    \begin{ex}
        Trong phản ứng $A + B \rightarrow C + 2D$, nồng độ chất A giảm từ 1.1M xuống 0.5M trong 10 phút. Tốc độ trung bình của phản ứng là:
         \choice
        {0.04 M/phút}
        {0.05 M/phút}
        {\True 0.06 M/phút}
        {0.08 M/phút}
        \loigiai{
            \[
                v_{tb} = \frac{1.1 - 0.5}{10} = 0.06 \text{ M/phút}
            \]
        }
    \end{ex}
     %%%%%============EX_20================%%%%%%
    \begin{ex}
        Trong phản ứng $A + B \rightarrow C + D$, nồng độ chất C tăng từ 0.1M lên 0.7M trong 12 phút. Tốc độ trung bình của phản ứng là:
        \choice
        {0.03 M/phút}
        {0.04 M/phút}
        {\True 0.05 M/phút}
        {0.06 M/phút}
        \loigiai{
            \[
                v_{tb} = \frac{0.7 - 0.1}{12} = 0.05 \text{ M/phút}
            \]
        }
    \end{ex}
\Closesolutionfile{ans}
\Closesolutionfile{ansex}
%\bangdapan{Ans-C06B01_Dang1}

%%%%%%%%%%%%%%%Trắc nghiệm đúng sai%%%%%%%%%%%%%%%%%%%%%%%%
\phan{Bài tập trắc nghiệm Đúng Sai}
%%%=============SOẠN EXTF===============%%%
\Opensolutionfile{ansex}[Ans/LGTF-C06B01_Dang1]
\Opensolutionfile{ansbook}[Ansbook/AnsTF-C06B01_Dang1]
\Opensolutionfile{ans}[Ans/Tempt-C06B01_Dang1]
    %%%%%============EXTF_1================%%%%%%
    \begin{ex}
        Tốc độ phản ứng là đại lượng đặc trưng cho sự biến thiên nồng độ của chất phản ứng theo thời gian.
        \choiceTF
        {\True Đúng}
        {\True Đúng}
        {\True Đúng}
        {\True Đúng}
        \loigiai{
            \begin{itemchoice}[T1,T2,T3,T4]
                \itemch Tốc độ phản ứng là đại lượng đặc trưng cho sự biến thiên nồng độ của chất phản ứng hoặc sản phẩm theo thời gian.
                \itemch Tốc độ phản ứng là đại lượng đặc trưng cho sự biến thiên nồng độ của chất phản ứng hoặc sản phẩm theo thời gian.
                \itemch Tốc độ phản ứng là đại lượng đặc trưng cho sự biến thiên nồng độ của chất phản ứng hoặc sản phẩm theo thời gian.
                \itemch Tốc độ phản ứng là đại lượng đặc trưng cho sự biến thiên nồng độ của chất phản ứng hoặc sản phẩm theo thời gian.
            \end{itemchoice}
        }
    \end{ex}
    %%%%%============EXTF_2================%%%%%%
    \begin{ex}
        Tốc độ trung bình của phản ứng luôn không đổi trong suốt quá trình phản ứng.
        \choiceTF
        {Sai}
        {\True Sai}
        {Sai}
        {Sai}
        \loigiai{
            \begin{itemchoice}[F1,T2,F3,F4]
                \itemch Tốc độ trung bình của phản ứng thường thay đổi trong suốt quá trình phản ứng.
                \itemch Tốc độ trung bình của phản ứng thường thay đổi trong suốt quá trình phản ứng.
                 \itemch Tốc độ trung bình của phản ứng thường thay đổi trong suốt quá trình phản ứng.
                \itemch Tốc độ trung bình của phản ứng thường thay đổi trong suốt quá trình phản ứng.
            \end{itemchoice}
        }
    \end{ex}
     %%%%%============EXTF_3================%%%%%%
    \begin{ex}
        Tốc độ phản ứng chỉ phụ thuộc vào nồng độ của các chất phản ứng.
        \choiceTF
        {Sai}
        {\True Sai}
        {Sai}
        {Sai}
        \loigiai{
            \begin{itemchoice}[F1,T2,F3,F4]
                \itemch Tốc độ phản ứng còn phụ thuộc vào nhiều yếu tố khác như nhiệt độ, áp suất, diện tích bề mặt, chất xúc tác.
                \itemch Tốc độ phản ứng còn phụ thuộc vào nhiều yếu tố khác như nhiệt độ, áp suất, diện tích bề mặt, chất xúc tác.
                 \itemch Tốc độ phản ứng còn phụ thuộc vào nhiều yếu tố khác như nhiệt độ, áp suất, diện tích bề mặt, chất xúc tác.
                \itemch Tốc độ phản ứng còn phụ thuộc vào nhiều yếu tố khác như nhiệt độ, áp suất, diện tích bề mặt, chất xúc tác.
            \end{itemchoice}
        }
    \end{ex}
     %%%%%============EXTF_4================%%%%%%
    \begin{ex}
        Tốc độ phản ứng luôn tăng khi tăng nồng độ các chất phản ứng.
        \choiceTF
        {\True Đúng}
        {\True Đúng}
         {\True Đúng}
        {\True Đúng}
        \loigiai{
            \begin{itemchoice}[T1,T2,T3,T4]
                \itemch Tốc độ phản ứng thường tăng khi tăng nồng độ các chất phản ứng.
                \itemch Tốc độ phản ứng thường tăng khi tăng nồng độ các chất phản ứng.
                \itemch Tốc độ phản ứng thường tăng khi tăng nồng độ các chất phản ứng.
                \itemch Tốc độ phản ứng thường tăng khi tăng nồng độ các chất phản ứng.
            \end{itemchoice}
        }
    \end{ex}
     %%%%%============EXTF_5================%%%%%%
    \begin{ex}
        Tốc độ phản ứng luôn tăng khi tăng nhiệt độ.
        \choiceTF
        {\True Đúng}
        {\True Đúng}
         {\True Đúng}
        {\True Đúng}
        \loigiai{
            \begin{itemchoice}[T1,T2,T3,T4]
                \itemch Tốc độ phản ứng thường tăng khi tăng nhiệt độ.
                \itemch Tốc độ phản ứng thường tăng khi tăng nhiệt độ.
                \itemch Tốc độ phản ứng thường tăng khi tăng nhiệt độ.
                \itemch Tốc độ phản ứng thường tăng khi tăng nhiệt độ.
            \end{itemchoice}
        }
    \end{ex}
     %%%%%============EXTF_6================%%%%%%
    \begin{ex}
        Tốc độ phản ứng luôn giảm khi có chất xúc tác.
        \choiceTF
        {Sai}
        {\True Sai}
        {Sai}
        {Sai}
        \loigiai{
            \begin{itemchoice}[F1,T2,F3,F4]
                \itemch Chất xúc tác làm tăng tốc độ phản ứng.
                \itemch Chất xúc tác làm tăng tốc độ phản ứng.
                \itemch Chất xúc tác làm tăng tốc độ phản ứng.
                \itemch Chất xúc tác làm tăng tốc độ phản ứng.
            \end{itemchoice}
        }
    \end{ex}
\Closesolutionfile{ans}
\Closesolutionfile{ansbook}
\Closesolutionfile{ansex}
\end{document}